% ------------------------------------------------------------------------------
% Resumo (NBR 6028: parágrafo único, 150 a 500 palavras, verbos na voz ativa)
% ------------------------------------------------------------------------------

\begin{resumo}

    O Programa Pé-de-Meia, instituído pela Lei nº 14.818/2024, é um incentivo financeiro-educacional na modalidade de poupança destinado a estudantes do ensino médio público pertencentes a famílias inscritas no Cadastro Único, com o objetivo de promover a permanência e a conclusão escolar.
    Este artigo realiza uma avaliação observacional, em escopo nacional, da associação entre a implantação do programa e os indicadores educacionais do ensino médio.
    Construiu-se um painel balanceado de 324 observações (27 Unidades da Federação $\times$ redes estadual e privada $\times$ seis anos entre 2019 e 2025), integrando microdados de pagamentos da Controladoria-Geral da União, taxas de rendimento do Censo Escolar e resultados do SAEB/IDEB do INEP.
    A execução financeira de 2024 foi auditada a partir das nove folhas mensais publicadas, totalizando R\$ 5,59 bilhões em 27,9 milhões de parcelas, e identificou-se a retificação de 1,79 milhão de registros duplicados na folha de março de 2024.
    Estimaram-se 21 modelos de Diferenças em Diferenças com efeitos fixos de UF $\times$ rede e de ano e erros-padrão agrupados por UF com distribuição $t$ (26 graus de liberdade).
    No período pós-programa (2024--2025), a rede estadual apresentou, em relação à privada, redução diferencial de 2,07 pontos percentuais na taxa de abandono e aumento de 4,28 p.p. na taxa de aprovação ($p < 0{,}001$).
    Na janela 2023 \textit{versus} 2025, o IDEB estadual avançou 0,18 ponto ($p = 0{,}0115$), impulsionado sobretudo pelo indicador de rendimento, sem ganho significativo na proficiência padronizada ($p = 0{,}1830$).
    Contudo, os testes formais rejeitam a hipótese de tendências paralelas no período pré-programa, o que, somado à multiplicidade de testes e à agregação ecológica, impede a atribuição causal.
    Os resultados devem, portanto, ser interpretados como associações descritivas, e o estudo disponibiliza um \textit{pipeline} reprodutível para avaliações futuras com microdados individuais.

    \par\vspace{\baselineskip}

    \noindent\textbf{Palavras-chave}: Pé-de-Meia. Diferenças em Diferenças. Abandono escolar. IDEB. Ciência de Dados.
\end{resumo}
